% Ad Atticum 12.36 --- headnote
% ad-atticum-12-36, 3 May 45 BC, Astura

\textit{Cicero to Atticus, written from Astura on the
fifth day before the Nones of May 709 AUC --- 3 May
45 BC (the manuscript dateline: \textit{Scr.\ Asturae
v Non.\ Mai.\ a.\ 709 (45)}). Cicero is back at
Astura after the brief journey to Sicca's
suburbanum and Rome. The first section is the most
explicit statement in the whole sequence of what the
\textit{fanum} is for: ``I want a shrine made, and
this cannot be torn from me.'' He is bent on
escaping any resemblance to a tomb, both to avoid the
funerary law's penalty and --- the deeper reason ---
to attain for Tullia an apotheosis,
\textit{[Greek: apothe\=osin]}: a deification in
the Greek manner.}

\textit{The aside on placing the shrine within the
villa rather than in open ground is among the most
human moments of the sequence: Cicero fears
\textit{commutationes dominorum}, the changes of
owners that would leave the shrine unprotected. In a
field, posterity will hold it religiously; in a
private house, it depends on the next purchaser. He
confesses these to be \textit{ineptiae} ---
``foolish notions,'' although the word in Latin
also has a sense of misplaced fussiness --- and asks
Atticus to send him the text of the funerary law to
study for evasions. The second section turns to
Brutus' refusal to visit the Cumanum (a small
\textit{obiurgatio} for Atticus to transmit) and to
the architect Cluatius, who is to be encouraged in
the work whatever site is finally chosen. The closing
``you yourself, perhaps, will be at the villa
tomorrow'' is the kind of plain practical note that
keeps appearing in the middle of the metaphysics.}
